\parttitle{VIII}{Design Principles and Research Directions}

\section{Design principles}
\label{sec:design-principles}

Ten principles condense Parts~I--VII. Each is stated as an imperative with
its justification and the place in this document where it was earned.

\emph{D1 --- Name your contraction.} Every working reservoir has an explicit
mechanism that erases initial conditions and the remote past: a partial
trace, a reset, engineered dissipation, a window truncation. State which one
your design uses, verify it numerically (trajectory convergence from
different $\rho_0$; PTM spectral radius), and check it contracts toward an
\emph{input-sensitive} region rather than the maximally mixed state
(coherence influx, Section~\ref{sec:quantum-esp}). Architectures that
inherit stability from unexamined device noise are not reproducible across
calibration cycles (Sections~\ref{sec:arch-dissipative}, \ref{sec:hw-ibm}).

\emph{D2 --- Sweep the operating point.} The injection interval and dynamical
regime are not defaults but the reservoir's most sensitive dials; the
temporal-edge sweep of Section~\ref{sec:temporal-edge} moved NARMA-10 error
by a factor of nearly four at fixed architecture. Re-sweep when the task,
input scaling, or noise level changes; on hardware the sweep reuses circuits
and is cheap (Figure~\ref{fig:temporal-edge};
\citealp{kobayashi2026edge,martinez2021dynamical}).

\emph{D3 --- Budget features, not qubits.} The read-out sees $K$ measured,
well-conditioned directions, not $4^n-1$ abstract ones. Track effective rank
and memory capacity of the actual feature matrix; enrich observables
($X$, $ZZ$) and multiplex before adding qubits
(Sections~\ref{sec:extraction}, \ref{sec:cirq}; Table~\ref{tab:cirq-sweep}).

\emph{D4 --- Treat the encoding scale as physics.} Rotation encodings place
the model's approximation power on a frequency comb set by the scale
$\gamma$ \citep{schuld2021effect}; the wrong scale wastes the register on
harmonics the task lacks. The single largest error reduction in Part~IV came
from halving $\gamma$, and the cheapest gradient in the document was the one
tuning it (Sections~\ref{sec:cirq}, \ref{sec:pennylane}).

\emph{D5 --- Anchors before experiments.} Build the validation suite first:
identities with known answers (injection rebuild, unitarity, ridge recovery,
NMSE of the mean $=1$, shot-noise scaling). Every later anomaly is then
attributable in minutes rather than days
(Section~\ref{sec:validation-suite}).

\emph{D6 --- Shots are the currency.} Measure the exact ceiling before any
sampled result; report $S$ next to every number; budget mitigation
multipliers explicitly; and read \citet{hu2023tackling} and
\citet{ahmed2025optimal} before a device campaign. A method that needs
$10^5$ shots per feature to approach its own ceiling has stated its cost
(Sections~\ref{sec:shot-practice}, \ref{sec:shot-budget}).

\emph{D7 --- The baseline battery is part of the model.} Mean, persistence,
linear lags, size-matched ESN, Haar control; chronological splits; the
physically stationary target ($k_t$, anomaly, residual). Part~VII's solar
row exists to show how publishable a meaningless result can look when the
battery is absent (Sections~\ref{sec:baselines}, \ref{sec:task-demos}).

\emph{D8 --- Default to the hybrid posture, with ablations.} Concatenate
quantum features with the cheap classical ones and report quantum-only and
classical-only columns; the difference is the measured quantum
contribution. Part~IV's noisy hybrid \emph{underperforming} its classical
ablation is exactly the kind of result this posture is designed to surface
(Sections~\ref{sec:arch-hybrid}, \ref{sec:qiskit-results}).

\emph{D9 --- Respect concentration and symmetry at scale.} Monitor feature
variance and effective rank as $n$, depth, or noise grow; expect exponential
signal collapse in scrambling or noise-dominated regimes and design against
it (shallow steps, local observables, spatial multiplexing); resolve
symmetry sectors before trusting chaos diagnostics or observable choices
(Section~\ref{sec:concentration};
\citealp{sannia2025exponential,xiong2025role,mcclean2018barren}).

\emph{D10 --- Scope every claim.} Theory, simulation, or hardware; surrogate
or real; the phrase ``quantum advantage'' only behind a controlled scaling
comparison that includes cost. The defensible current claim for QRC in this
domain is parameter-parity plus engineering fit, and this document says no
more anywhere (frontmatter; Section~\ref{sec:arch-comparison}).

\section{Detecting real quantum utility}
\label{sec:detecting-utility}

A benchmark that could in principle certify a quantum benefit has five
properties, and their absence diagnoses most overclaims. \emph{Matched
read-out capacity}: classical baselines get the same number of trained
parameters and comparable tuning effort. \emph{Structure controls}: a
Haar-random quantum reservoir of the same size separates ``tuned dynamics''
from ``being quantum'' (Figure~\ref{fig:memory-capacity}); an efficiently
simulable quantum control (Gaussian or Clifford reservoir) separates
``quantum hardware'' from ``quantum formalism''
\citep{nokkala2021gaussian,gotting2023exploring,palacios2024role}.
\emph{Scaling curves, not points}: performance versus size, shots, and data
on both sides, since crossings, not gaps at one configuration, are the
object of interest. \emph{Cost on the axis}: wall-clock, shots, and energy;
Section~\ref{sec:shot-budget}'s arithmetic belongs in every hardware paper.
\emph{Pre-registered success criteria}: the target, metric, split, and
baselines fixed before the test set is touched --- the cheapest and most
neglected of the five.

\section{Research directions}
\label{sec:research-directions}

Ten directions follow, ordered roughly by increasing ambition. Each gives
motivation, background, a concrete route, difficulty, data, metrics, and
publication potential. All are scoped so that a small group with the shipped
code and modest cloud access can start immediately.

\subsection*{R1 --- A temporal-edge atlas across architectures}
\emph{Motivation:} D2 is currently a per-study ritual; a systematic map
would make it a lookup. \emph{Background:}
\citet{kobayashi2026edge} establish the two-edge (parametric, temporal)
structure for SYK-class reservoirs; Part~VII reproduced the temporal edge in
a 5-qubit Ising register. \emph{Route:} sweep $(J\Delta t$, disorder,
task-memory-length$)$ across the Part~III templates (Ising, chain,
brickwork circuit) at $n\le 12$ exact, correlating optima with level-spacing
ratio and PTM gap; deliver a design chart and the sweep harness.
\emph{Difficulty:} low--medium (compute-bound, conceptually clean).
\emph{Data:} NARMA family, solar $k_t$, ENSO. \emph{Metrics:} NMSE surfaces;
optimum-location predictors. \emph{Potential:} strong methods paper
(PRApplied/PRResearch class); immediately useful tooling.

\subsection*{R2 --- Closing the recurrent-memory gap on ENSO-class tasks}
\emph{Motivation:} Part~VII's sharpest result --- recurrence worth 20 NMSE
points over windowed maps, tuned ESN still ahead of the small register.
\emph{Background:} memory-augmented hybrids \citep{settino2025memory},
feedback reservoirs \citep{kobayashi2024feedback}, NISQRC's
coherence-independent memory \citep{hu2024overcoming}.
\emph{Route:} implement feedback and memory-augmented variants on the exact
simulator; identify which mechanism closes the gap to the matched ESN at
fixed feature count; then a small hardware run (mid-circuit reset, Part~IV
Part-B shape) on the real ENSO series. \emph{Difficulty:} medium--high.
\emph{Data:} ENSO (real); Ni\~no 3.4 from reanalysis for length.
\emph{Metrics:} NMSE and warm-event skill versus matched ESN.
\emph{Potential:} first honest hardware benchmark on a real climate index;
high.

\subsection*{R3 --- Shot-optimal read-outs for forecasting}
\emph{Motivation:} Section~\ref{sec:shot-budget}'s arithmetic is the
adoption barrier. \emph{Background:} eigentask construction
\citep{hu2023tackling}; shot-optimal training \citep{ahmed2025optimal}.
\emph{Route:} apply eigentask-weighted observable selection to the solar and
load pipelines; produce NMSE-versus-total-shot-budget frontiers against
naive uniform allocation. \emph{Difficulty:} medium.
\emph{Data:} surrogate then NSRDB/GEFCom. \emph{Metrics:} error at fixed
budget; budget at fixed error. \emph{Potential:} the paper every hardware
group cites in its methods; solid.

\subsection*{R4 --- Event skill: ramps and extremes}
\emph{Motivation:} Part~VII shows averages are ties; operational value
lives in tails. \emph{Background:} extreme-event reservoir prediction
\citep{ahmed2024prediction}; energy community's event metrics
\citep{hong2016gefcom,inman2013solar}. \emph{Route:} define ramp events on
SURFRAD minute data and warm events on ENSO; train quantile read-outs
(Section~\ref{sec:uncertainty}); compare event-conditional skill (POD/FAR,
tail pinball) across the reservoir battery. \emph{Difficulty:} medium
(metrics discipline is the work). \emph{Data:} SURFRAD (real), ENSO (real).
\emph{Metrics:} event scores, reliability diagrams. \emph{Potential:} high
--- the first tail-focused QRC study in the domain would define the
evaluation standard.

\subsection*{R5 --- Noise as a characterised resource, on hardware}
\emph{Motivation:} D1/D5 versus the reality that device noise drifts.
\emph{Background:} natural reservoirs \citep{suzuki2022natural};
dissipation-as-resource \citep{sannia2024dissipation}; coherence-influx
criterion \citep{kobayashi2024coherence}. \emph{Route:} on one
superconducting backend, characterise the same reservoir circuit across
calibration cycles; correlate task performance with measured
$T_1/T_2$/read-out drift; test the amplitude-damping-helps /
dephasing-hurts prediction by Trotter-inserting engineered channels in
simulation first. \emph{Difficulty:} medium (patience-bound).
\emph{Data:} NARMA + solar surrogate. \emph{Metrics:} performance-drift
curves; influx diagnostics. \emph{Potential:} a reproducibility protocol
the field lacks; solid.

\subsection*{R6 --- The hybrid load benchmark, pre-registered}
\emph{Motivation:} Section~\ref{sec:pipeline-load} is specified but unrun;
it is the cleanest first real-data project. \emph{Background:} GEFCom
practice \citep{hong2016gefcom,hong2016load}; QRC load studies at few-qubit
scale \citep{odo2026quantum}. \emph{Route:} execute the pipeline exactly as
specified (calendar model, residual target, hybrid read-out, ablations,
pinball quantiles) on GEFCom2014 and one ENTSO-E zone; pre-register the
success criterion. \emph{Difficulty:} low--medium.
\emph{Data:} GEFCom2014, ENTSO-E/OPSD (all real).
\emph{Metrics:} pinball loss versus calendar-plus-linear; ablation deltas.
\emph{Potential:} modest novelty, high credibility value --- the reference
others must match methodologically.

\subsection*{R7 --- Solar beyond persistence: multi-horizon, regime-aware}
\emph{Motivation:} the one-step tie (Table~\ref{tab:results-master}) says
move the goalposts to where value exists. \emph{Background:} solar
benchmarking practice \citep{yang2018history,inman2013solar}; hybrid
model-plus-reservoir posture \citep{pathak2018hybrid}.
\emph{Route:} 2--6\,h horizons on NSRDB/SURFRAD $k_t$ with
regime-conditional evaluation; test whether reservoir memory helps most at
regime transitions; optionally fuse a cheap clear-sky/NWP prior.
\emph{Difficulty:} medium. \emph{Data:} NSRDB, SURFRAD (real).
\emph{Metrics:} NMSE by horizon and regime; ramp scores (with R4).
\emph{Potential:} directly relevant to the application community; good
cross-venue paper.

\subsection*{R8 --- Empirical concentration scaling for QRC}
\emph{Motivation:} D9 currently cites theory; design needs measured
curves. \emph{Background:} exponential concentration and symmetry analysis
\citep{sannia2025exponential,xiong2025role}; eigentask noise ceiling
\citep{hu2023tackling}. \emph{Route:} exact simulation to $n=12$--$14$:
feature variance, effective rank, and task NMSE versus $n$, depth, noise
rate, and symmetry sector; fit decay constants; identify the multiplexing
crossover where two $n$-qubit registers beat one $2n$-qubit register at
fixed shots. \emph{Difficulty:} medium--high (careful numerics).
\emph{Data:} NARMA suffices. \emph{Metrics:} scaling exponents; crossover
maps. \emph{Potential:} the quantitative backbone for honest scaling
claims; strong.

\subsection*{R9 --- An analog-platform energy demonstration}
\emph{Motivation:} the largest registers are analog
(Section~\ref{sec:hw-neutral}); the domain lacks any energy-data result on
them. \emph{Background:} 108-qubit Rydberg QRC \citep{kornjaca2024large};
global-detuning encoding's size-independent memory.
\emph{Route:} map a short-window load or ramp \emph{classification} task
(matching the platform's restart/ensemble operation) onto global-detuning
waveforms in emulation; then a partner-hardware run; report against the
full battery including a matched classical kernel.
\emph{Difficulty:} high (platform access and encoding design).
\emph{Data:} GEFCom/SURFRAD-derived labels. \emph{Metrics:} accuracy at
matched trained parameters; shots and wall-clock on the axis.
\emph{Potential:} headline-adjacent if honest; must be pre-registered to
stay so.

\subsection*{R10 --- Clean-physics ESP and edge studies on trapped ions}
\emph{Motivation:} Part~V flags the absence of ion-platform QRC benchmarks
despite the best per-step fidelities and native all-to-all coupling.
\emph{Background:} QCCD mid-circuit primitives \citep{moses2023race};
dynamical-phase reservoir physics
\citep{martinez2021dynamical,kobayashi2026edge}. \emph{Route:} small-$T$,
high-fidelity implementations of the Eq.~\eqref{eq:fn-hamiltonian}
reservoir with mid-circuit reset; measure trajectory-convergence rates
(ESP directly, not via proxy) and the temporal edge on-device; compare with
exact simulation at identical seeds. \emph{Difficulty:} high (access;
shot-rate limits set $T$). \emph{Data:} NARMA-2/10 at small $T$.
\emph{Metrics:} measured contraction rates; edge location on-device versus
simulation. \emph{Potential:} the cleanest hardware--theory comparison
available to the field; high scientific value independent of any advantage
question.

\section{Closing statement}
\label{sec:closing}

This document set out to be the reference its own Part~VIII demands:
theory stated with its assumptions, implementations that run, data
characterised before being forecast, baselines that would embarrass a weak
claim, and research directions scoped to be started rather than admired.
The honest summary of the field it surveys is neither triumphal nor
dismissive. Small quantum registers, operated at their dynamical sweet
spot, demonstrably match classical read-outs several times their size on
memory-nonlinearity benchmarks, and the engineering paths around their two
structural taxes --- measurement statistics and the recurrence gap --- are
active and credible. Whether that matures into practical value for
renewable-energy and climate forecasting is exactly the kind of question
that is settled by pre-registered benchmarks on real data, run by people
who would publish either answer. The code, data discipline, and directions
above are offered to make that settling faster.
